\section{Marco teorico del modelo M/M/1}

\subsection{Supuestos del modelo}

El modelo M/M/1 representa un sistema de espera con una sola estacion de servicio. La primera M indica que los tiempos entre arribos siguen una distribucion exponencial, equivalente a un proceso de Poisson para la cantidad de arribos. La segunda M indica que los tiempos de servicio tambien son exponenciales. El numero 1 indica que existe un unico servidor.

En este trabajo se asume disciplina FIFO, es decir, los clientes son atendidos en el mismo orden en el que llegan. Las tasas se expresan en clientes por hora:

\begin{itemize}
    \item $\lambda$: tasa media de arribos.
    \item $\mu$: tasa media de servicio.
    \item $\rho = \lambda / \mu$: factor de carga o utilizacion teorica del sistema con cola infinita.
\end{itemize}

\subsection{Cola infinita}

Para una cola infinita, el sistema solo tiene regimen estacionario si:

\begin{equation}
    \rho < 1.
\end{equation}

Cuando esta condicion se cumple, las medidas teoricas usadas para validar la simulacion son:

\begin{align}
    L &= \frac{\rho}{1-\rho}, \\
    L_q &= \frac{\rho^2}{1-\rho}, \\
    W &= \frac{1}{\mu-\lambda}, \\
    W_q &= \frac{\lambda}{\mu(\mu-\lambda)}, \\
    U &= \rho.
\end{align}

Donde $L$ es el numero promedio de clientes en el sistema, $L_q$ es el numero promedio de clientes en cola, $W$ es el tiempo promedio en el sistema, $W_q$ es el tiempo promedio de espera en cola y $U$ es la utilizacion del servidor.

Para la distribucion de longitud de cola se usa la relacion entre el numero de clientes en sistema $N$ y el numero de clientes esperando $Q$. Si $Q=0$, el sistema puede estar vacio o tener un cliente en servicio. Por lo tanto:

\begin{align}
    P(Q=0) &= P(N=0) + P(N=1), \\
    P(Q=q) &= P(N=q+1), \quad q \geq 1.
\end{align}

Como en M/M/1 estable $P(N=n)=(1-\rho)\rho^n$, queda:

\begin{align}
    P(Q=0) &= (1-\rho)(1+\rho), \\
    P(Q=q) &= (1-\rho)\rho^{q+1}, \quad q \geq 1.
\end{align}

Si $\rho \geq 1$, la cola infinita no tiene valores estacionarios finitos. En esos casos la simulacion por horizonte finito puede generar metricas observadas, pero no corresponde compararlas contra las formulas estacionarias anteriores.

\subsection{Cola finita}

Para estudiar denegacion de servicio se usa una cola con capacidad finita. En el codigo y en las tablas, $B$ representa la cantidad maxima de clientes que pueden esperar en cola. Como ademas puede haber un cliente en servicio, la capacidad total del sistema es:

\begin{equation}
    K = B + 1.
\end{equation}

El estado $n$ representa la cantidad de clientes dentro del sistema, con $0 \leq n \leq K$. Las probabilidades estacionarias para M/M/1/K son:

\begin{equation}
    P_n =
    \begin{cases}
        \frac{(1-\rho)\rho^n}{1-\rho^{K+1}}, & \rho \neq 1, \\
        \frac{1}{K+1}, & \rho = 1.
    \end{cases}
\end{equation}

La probabilidad de denegacion es la probabilidad de que el sistema este lleno al momento de un arribo:

\begin{equation}
    P_{den} = P_K.
\end{equation}

La tasa efectiva de clientes aceptados queda:

\begin{equation}
    \lambda_{\mathrm{eff}} = \lambda(1-P_{den}).
\end{equation}

Con estas probabilidades, las medidas teoricas se calculan como:

\begin{align}
    L &= \sum_{n=0}^{K} n P_n, \\
    L_q &= \sum_{n=0}^{K} \max(n-1,0) P_n, \\
    W &= \frac{L}{\lambda_{\mathrm{eff}}}, \\
    W_q &= \frac{L_q}{\lambda_{\mathrm{eff}}}, \\
    U &= 1-P_0.
\end{align}

A diferencia de la cola infinita, el modelo con capacidad finita tiene regimen estacionario tambien para $\rho \geq 1$, porque los arribos excedentes se rechazan cuando el sistema esta lleno.
